\section{Comparison with gold antennas on ZnO}
\label{sec:zno}
Imasaka \textit{et al.}~\cite{imasaka2018} measured harmonics and band-gap fluorescence from gold
bars (length 570~nm, width 310~nm, thickness 100~nm) on $a$-cut ZnO driven by 100-fs pulses at
2~$\mu$m, incident from the substrate side and polarized along the bars, at intensities between
0.012 and 0.1~TW/cm$^2$~\cite{imasaka2019,ashihara2020}. Their simulation gives an enhancement of
the in-plane field of about 9 (intensity 80) at the hot spots, 2~nm below the surface. With the
bars along the $m$ axis only odd harmonics appear; along the polar $c$ axis even harmonics appear as
in the bare crystal. We compare with the intensity dependence of the fifth, seventh, and ninth
harmonics and of the fluorescence, which was reported for the $c$ orientation. From the published
double-logarithmic plot (Fig.~8 of Ref.~\cite{ashihara2020}) we read power-law exponents of
3.4, 3.2, 3.0, and 2.7, respectively, with an uncertainty of about 0.3; the multiphoton laws would
give 5, 7, and 9 for the harmonics.

\subsection{Near field of the bar}
Near the contact edge the field is quasi-static. We solve
$\nabla\cdot(\varepsilon\nabla\phi)=0$ by finite differences on a graded grid (smallest spacing
0.25~nm) in the cross-section along the bar axis, for a perfectly conducting bar of length 570~nm
and height 100~nm on a half-space with $\varepsilon=3.7$ in a uniform field along the bar.
Figure~\ref{figzno}(a) shows the result. The field is singular at the contact line; averaged over
depths of 0.5 to 2~nm, it decays to one half of its value at the edge within 5~nm on either side,
and below the edge it decays with depth approximately as $z^{-1/2}$. The calculation does not
contain the antenna resonance. We therefore use it for the shape of the near field only and fix
the magnitude by a single number, the enhancement $\eta$ of the peak of $|E_x|$ at 2~nm depth
relative to the incident amplitude inside the substrate, $0.60$~V/nm at 0.1~TW/cm$^2$.

\subsection{Chain model with the computed depth profile}
We repeated the field scan of the chain model with the depth dependence of the field below the edge
(1~nm inside the contact) in place of the exponential with $\ell=8$~nm. Table~\ref{tab:prof} shows
that the rectification is somewhat stronger than for the exponential profile: at $F_0=1.1$~V/nm the
forward half-cycle is reduced by 24\%, the reverse one enhanced by 40\%, and
$F_{\rm dc}=0.28F_0$. The table also lists the density of holes left in the crystal after the
pulse, obtained from the depletion of the initially occupied eigenstates weighted with their
probability beyond the second crystal layer. We use it as a proxy for the band-gap fluorescence,
which requires holes in the $n$-type crystal. The contact raises the hole density at fields below 2~V/nm.

\begin{table}[h]
\caption{Chain model with the depth profile of the near field below the edge of the bar. Other
parameters as in Table~\ref{tab:field}.}
\label{tab:prof}
\begin{tabular}{lcccccccccc}
\hline\hline
$D_0$ & $F_0$ & $F_-^{\rm ins}$ & $F_-$ & $F_+$ & $F_{\rm dc}$ & $W_{\rm contact}$ & $W_{\rm insulated}$ & \multicolumn{2}{c}{holes (cm$^{-2}$)} & $S_2$ ratio \\
\multicolumn{6}{c}{(V/nm)} & \multicolumn{2}{c}{(eV/nm$^2$)} & contact & insulated & \\
\hline
1.3 & 0.31 & 0.28 & 0.25 & 0.30 & 0.01 & 0.0072 & $3.4\times10^{-5}$ & $7.6\times10^{9}$ & $4.5\times10^{8}$ & 491 \\
1.9 & 0.45 & 0.41 & 0.36 & 0.47 & 0.04 & 0.0316 & $3.4\times10^{-5}$ & $7.1\times10^{10}$ & $4.8\times10^{8}$ & 425 \\
2.6 & 0.62 & 0.56 & 0.48 & 0.71 & 0.10 & 0.0896 & $6.0\times10^{-5}$ & $1.7\times10^{11}$ & $1.1\times10^{9}$ & 291 \\
3.2 & 0.75 & 0.69 & 0.57 & 0.94 & 0.15 & 0.1200 & $2.5\times10^{-4}$ & $3.9\times10^{11}$ & $5.6\times10^{9}$ & 136 \\
4.0 & 0.93 & 0.86 & 0.68 & 1.25 & 0.23 & 0.2473 & $1.9\times10^{-3}$ & $1.1\times10^{12}$ & $4.2\times10^{10}$ & 41 \\
4.8 & 1.10 & 1.03 & 0.78 & 1.54 & 0.31 & 0.4822 & $9.4\times10^{-3}$ & $3.0\times10^{12}$ & $2.0\times10^{11}$ & 26 \\
5.6 & 1.26 & 1.20 & 0.93 & 1.78 & 0.32 & 0.6661 & $3.2\times10^{-2}$ & $6.4\times10^{12}$ & $6.7\times10^{11}$ & 13 \\
6.5 & 1.45 & 1.41 & 1.23 & 1.99 & 0.34 & 1.2086 & $1.1\times10^{-1}$ & $1.4\times10^{13}$ & $2.2\times10^{12}$ & 6 \\
8.0 & 1.73 & 1.83 & 1.69 & 2.33 & 0.38 & 3.3076 & $4.7\times10^{-1}$ & $4.2\times10^{13}$ & $9.1\times10^{12}$ & 5 \\
10.0 & 2.01 & 2.40 & 2.34 & 3.01 & 0.41 & 9.6042 & $2.2\times10^{0}$ & $1.2\times10^{14}$ & $4.0\times10^{13}$ & 9 \\
13.0 & 3.14 & 3.03 & 2.93 & 3.52 & 0.32 & 26.0142 & $1.2\times10^{1}$ & $3.5\times10^{14}$ & $2.1\times10^{14}$ & 10 \\
17.0 & 3.72 & 4.74 & 3.09 & 5.32 & 0.08 & 58.6693 & $5.1\times10^{1}$ & $8.7\times10^{14}$ & $8.5\times10^{14}$ & 7 \\
\hline\hline
\end{tabular}

\end{table}

\subsection{Hot-spot average}
The response per unit length of the edge is
\begin{equation}
X(I)=\int_{\rm metal}X_{\rm c}\big(|F_z(s)|\big)\,ds+\int_{\rm free}X_{\rm i}\big(|{\bf F}(s)|\big)\,ds ,
\end{equation}
where $s$ is the distance from the contact edge, $F_z$ and ${\bf F}$ are the fields of
Fig.~\ref{figzno}(a) averaged over depths of 0.5 to 2~nm and scaled with $\eta$ and the incident
amplitude, and $X_{\rm c}$ and $X_{\rm i}$ are the results of the contacted and the insulated chain
(harmonic yield, hole density, or absorbed energy per area), interpolated in the local field. The
integrals extend to the middle of the bar and 150~nm beyond its end. For a hypothetical insulated
antenna $X_{\rm c}$ is replaced by $X_{\rm i}$. Adding yields rather than amplitudes neglects the
phases of the harmonics, and the five-cycle pulse of the simulation is shorter than that of the
experiment; both affect magnitudes more than exponents.

Figure~\ref{figzno}(b) and Table~\ref{tab:zno} show the exponents over the measured range of
intensities as a function of $\eta$. They decrease with $\eta$ because a larger part of the hot spot
saturates. For $\eta=3$ to 4 the model reproduces the four measured exponents to within 0.5, except
for the ninth harmonic at $\eta=3$ and the seventh at $\eta=4$ (0.7 to 1.1). The harmonics do not
distinguish between contacted and insulated antennas. The fluorescence does: with the contact its
exponent is 2.7 to 2.9, as measured, and without it 3.3 to 4.0. For the simulated peak enhancement
$\eta=9$~\cite{imasaka2019} all exponents are near 2, below the measured ones, with or without the
contact; the peak field would then reach 5~V/nm at 0.1~TW/cm$^2$, where the model predicts an absorbed
energy of the order of 100~eV/nm$^2$ at the edge. We conclude that the measured intensity dependence is consistent with
the model for an effective enhancement of 3 to 4, and that within this description the fluorescence
exponent requires the contact. The effective enhancement is the one adjusted parameter.

Figure~\ref{figzno}(c) shows where the contact matters. For $\eta=4$ it raises the absorbed energy
by a factor of 27 and the second harmonic of a single end by 35 at 0.012~TW/cm$^2$; both factors
fall below 5 at 0.1~TW/cm$^2$, and they exceed 100 below 0.004~TW/cm$^2$. For $\eta=9$ the same
factors are reached at intensities five times lower. The contact thus dominates at the low end of
the intensities at which antenna-enhanced harmonics are detected and is masked by interband
excitation at the high end.

\begin{figure}[h]
\includegraphics[width=0.98\textwidth]{figS_zno.pdf}
\caption{(a) Quasi-static near field of a gold bar (570~nm $\times$ 100~nm) on ZnO around the
contact edge, averaged over depths of 0.5 to 2~nm and normalized to the peak of $|E_x|$ at 2~nm
depth. (b) Exponents of the intensity dependence of the hot-spot-averaged fluorescence proxy and
seventh harmonic between 0.012 and 0.1~TW/cm$^2$ versus the peak field enhancement, with
(circles) and without (squares) the contact; the bands mark the measured
values~\cite{ashihara2020}. (c) Ratio of the hot-spot-averaged absorbed energy and second-harmonic
yield with and without the contact for an enhancement of 4 (solid) and 9 (dashed); the shaded region
is the measured range.}
\label{figzno}
\end{figure}

\begin{table}[h]
\caption{Exponents of the intensity dependence between 0.012 and 0.1~TW/cm$^2$ for the contacted
antenna (in parentheses: insulated antenna) compared with the measured
ones~\cite{imasaka2018,ashihara2020}. The fluorescence is represented by the hole density.}
\label{tab:zno}
\begin{tabular}{llcccc}
\hline\hline
Near-field profile & Enhancement & 5th harmonic & 7th harmonic & 9th harmonic & Fluorescence \\
\hline
computed profile & 3 & 3.3 (3.5) & 3.1 (3.0) & 4.1 (3.6) & 2.9 (4.0) \\
computed profile & 4 & 3.0 (3.0) & 2.5 (2.6) & 3.0 (3.0) & 2.7 (3.3) \\
computed profile & 5 & 3.0 (2.9) & 2.6 (2.5) & 2.5 (2.9) & 2.4 (2.8) \\
computed profile & 9 & 2.4 (2.5) & 2.3 (2.1) & 2.3 (2.2) & 1.8 (2.0) \\
\hline
$\ell=8$ nm & 3 & 2.9 (3.0) & 2.7 (2.6) & 3.7 (3.7) & 3.2 (4.1) \\
$\ell=8$ nm & 4 & 2.7 (2.9) & 2.4 (2.4) & 2.2 (2.1) & 3.2 (3.7) \\
$\ell=8$ nm & 5 & 2.7 (2.9) & 2.3 (2.4) & 1.7 (1.6) & 2.8 (3.2) \\
$\ell=8$ nm & 9 & 2.5 (2.2) & 2.1 (2.0) & 1.8 (1.6) & 1.7 (1.8) \\
\hline
Measured & & 3.4 & 3.2 & 3.0 & 2.7 \\
Multiphoton law & & 5 & 7 & 9 & \\
\hline\hline
\end{tabular}

\end{table}
